\section{Comparison with published model structures}
\label{si:s3}

Every estimate of residual capacity is made within a model. This section asks which of our
conclusions still hold when the records are generated by the two published
pharmacokinetic\slash pharmacodynamic structures for mRNA-3927 instead of by our metabolite balance,
and analysed, as a trial analyst would, with the linear-clearance model. The ranking of sampling
schedules is the same under every structure tested, whereas the value of $\hat\theta$ changes
with the structure; one blood draw a few minutes after an infusion distinguishes the family of
structures to which the linear-clearance model belongs from the algebraic one. Supplementary
Table~\ref{stab:structures} collects the numbers. All records are at normal capacity
($\theta=3000$\um).

\subsection{The structures compared}
\label{si:s3:structures}

In the linear-clearance model the marker is the state of a flux balance: it relaxes towards the
level at which
the load is cleared by the enzyme, $(\theta+\Delta)\,m/(K_m+m)$, and by the non-enzymatic route,
with a time constant of about 0.8\,h (Section~\ref{si:s2:firsthour}). Between infusions it is in
quasi-steady state and obeys Eq.~\eqref{seq:qss}.
% src: results/numbers/design.tex:110; manuscript/main.tex:593-603,631-634

\emph{The algebraic structure} of Attarwala et al.~\cite{pkpd_translational} describes the
synthesis of PCC protein from plasma mRNA with an indirect-response model and an effect
compartment, and makes each marker a direct function of the enzyme, with no state of its own:
\begin{equation}
\label{seq:algebraic}
m \;=\; E_0+\mathrm{base}\Bigl(1-\frac{I_{\max}\,E}{\mathrm{IC}_{50}+E}\Bigr),
\end{equation}
where $E$ is the enzyme, $E_0$ a floor that the therapy cannot remove, base the part of the
marker it can suppress, $I_{\max}$ the maximal inhibition and $\mathrm{IC}_{50}$ the enzyme level
of half-maximal effect. Two of its constants transfer to our setting unchanged: $I_{\max}$, fixed
at 0.999 in the original fit, and the ratio $E_0/\mathrm{base}=0.0987/60.1$. Its
$\mathrm{IC}_{50}$ (37.5\,\textmu g/g) is a liver protein concentration, whereas $\theta$ is a flux
capacity, and the conversion between the two is not published, so $\mathrm{IC}_{50}$ on our
enzyme axis was calibrated as described below.
% src: audit/PHASE5_ROBUSTNESS.md:24-38,360-371,431-442,467-472; paper/claims.csv:42

\emph{The turnover structure} is the indirect-inhibition model that Gao et al.\ used for
3-hydroxy\-propionate and 2-methyl\-citrate in selecting the dose~\cite{gao2025acop}, a class also
named in the Methods of Baek et al.~\cite{mrna_moa}. In it the marker is a state whose
production the enzyme inhibits and whose loss is first order,
\begin{equation}
\label{seq:turnover}
\dot m \;=\; k_{\mathrm{in}}\Bigl(1-\frac{I_{\max}\,E}{\mathrm{IC}_{50}+E}\Bigr)-k_{\mathrm{out}}\,m .
\end{equation}
Neither source publishes its equations or parameter values, so we wrote down this standard form,
used the same $I_{\max}$, and calibrated $\mathrm{IC}_{50}$ and $k_{\mathrm{out}}$ to our operating
point.
% src: audit/PHASE5_ROBUSTNESS.md:68-90,401-422; audit/k2_structures.py:1-52; paper/claims.csv:42

In both structures we set $E=\theta+\Delta(t)$, with the dose train of Section~\ref{si:s2:setting}
($D=6000$\um\ every 14 days, half-life 3.5\,d), so the comparison concerns how the marker
responds to the enzyme rather than the pharmacokinetics of the mRNA. Neither published structure
contains a substrate load: the algebraic structure was fitted in mice on a fixed diet, and the
turnover structure is described only by name. Where the load enters is therefore our assumption.
In the turnover structure we took it as the marker's production rate, $k_{\mathrm{in}}=L+b$. In the
algebraic structure we scaled the whole marker by $(L+b)/\bar L$, with $\bar L=323$\um\ the
baseline load, so that at fixed enzyme the marker is proportional to the substrate supply;
Section~\ref{si:s3:load} examines the alternatives.
% src: audit/k2_structures.py:21-33,188-191 (PubB.steady scales by load/LB); audit/PHASE5_ROBUSTNESS.md:454-463

Each structure was calibrated against itself. The untreated marker was placed at our operating
level (10\uM\ at $\theta=3000$\um), and $\mathrm{IC}_{50}$ was chosen so that the marker falls by a
target fraction between the untreated enzyme and the steady-state peak of the treated enzyme, a
3.13-fold increase. Each target is met by two values of $\mathrm{IC}_{50}$, and both were kept. A
60\% fall stands for our assumed dose, which gives 66\% at the peak; 30\% and 20\% stand for a dose
calibrated to the mean falls in 3-hydroxypropionate reported in the trial, the largest of which is
29\%~\cite{trial_phase12}. For the 60\% target the two roots are $\mathrm{IC}_{50}=63.75$ and
1167\um\ in the algebraic structure and 22.93 and 1230\um\ in the turnover structure. The two roots
can imply very different marker memories: at the 30\% target, $k_{\mathrm{out}}$ of the turnover
structure is either 0.058\,d$^{-1}$ (a marker half-life of 12 days) or 25.8\,d$^{-1}$ (39 minutes).
The published inhibition shape also limits how far the marker can fall for this increase in
enzyme, at most 65.65\% in the turnover structure and 64.13\% in the algebraic structure with its
published floor, whatever $\mathrm{IC}_{50}$. In the linear-clearance model our assumed dose
gives 64.9\% at the first peak, beyond the algebraic structure's limit, and 66.4\% at the
steady-state peak, beyond both limits. This is an independent sign that the assumed dose is high
(Section~\ref{si:s2:record}).
% src: audit/k2_structures.py:120-167; audit/PHASE5_ROBUSTNESS.md:264-288,474-494,526-529; results/tables/tab_S_dosesens.tex:5

\subsection{Records and fits}
\label{si:s3:fits}

Records of 12 visits were generated on five schedules: at the dose interval from the start of
therapy, at the dose interval for a patient already on therapy, every 14.5 days, monthly and
phase-spread, with intake varying 0, 2 or 10\% from day to day and 8 intake records for each
setting. Each record was fitted with the linear-clearance model by variable projection, under the
priors of the main analysis (standard deviation 0.8 on $\ln k_e$, 5\% of the load on $b$; none on
$\theta$ or $K_m$). Noiseless records isolate the misfit due to structure; records carrying the
observation noise of the main analysis (5\% assay coefficient of variation, 5\% error in the
recorded load) show what an analyst would see. Misfit is reported as the root-mean-square residual in
units of the ordinate noise, $\sigma=22.0$\um. For a correct model fitted to 12 noisy draws with four
parameters, the expected value is $\sqrt{8/12}=0.816$. As a control, the linear-clearance model
fitted to its own records recovers $\theta$ exactly on every schedule (residual below $3\times10^{-13}\sigma$), except
at the dose interval for a patient already on therapy at a constant load, where the design is
exactly singular (Section~\ref{si:s2:aliasing}) and $\hat\theta$ goes to the estimator's lower
bound.
% src: audit/PHASE5_ROBUSTNESS.md:498-509,598-602,637-650; audit/k2_structures.py:433-520; results/numbers/design.tex:107

\begin{table}[htbp]
\centering\small\setlength{\tabcolsep}{4pt}
\caption{\label{stab:structures}Records generated by the published structures and analysed with
the linear-clearance model, at normal capacity, 12 visits per record. Algebraic: the $I_{\max}$
structure of Attarwala et al.~\cite{pkpd_translational}; turnover: the indirect-inhibition
structure of Gao et al.~\cite{gao2025acop}; calibration, the marker fall between the untreated
enzyme and the treated peak to which $\mathrm{IC}_{50}$ (\textmu M/d, on our enzyme axis) was
fitted. (a)~Noiseless records at a constant load: recovered $\hat\theta/\theta$ on each schedule
(dose int., at the dose interval from the start of therapy) and the residual of the monthly fit in
units of $\sigma$; 0 (bound), $\hat\theta$ at the estimator's lower bound; ranges cover both
$\mathrm{IC}_{50}$ roots. (b)~Records with the observation noise of the main analysis (5\% assay,
5\% load-record error), constant load, 60\% calibration: $\sdl{\hat\theta}$ as the fit reports
it, median over 8 noisy records; sing., singular information at every record.
(c)~Four ways of adding the load to the algebraic structure, and the structure as published for
reference, with $f=1-I_{\max}E/(\mathrm{IC}_{50}+E)$ and $u=(L+b)/\bar L$ ($\bar L$, the baseline
load), at a 60\% calibration ($\mathrm{IC}_{50}=63.75$) and 10\% day-to-day intake variation:
whether $\theta$ is identifiable in the generating structure at a varying recorded load; residual of the fit
(monthly, phase-spread) on noiseless records and with noise (expected 0.816 for a correct model);
$\hat\theta/\theta$ on noiseless monthly records with our priors and without them (medians over 8
records).}
\begin{tabular}{@{}llcccccc@{}}
\toprule
\multicolumn{8}{@{}l}{(a) recovered capacity and residual, noiseless records}\\
\midrule
generator & calibration & $\mathrm{IC}_{50}$ & dose int. & 14.5\,d & monthly & phase-spread & residual \\
\midrule
linear-clearance (control) & -- & -- & 1 & 1 & 1 & 1 & 0.000 \\
algebraic & 60\% & 63.75 & 2.03 & 1.53 & 1.34 & 1.32 & 0.180 \\
algebraic & 60\% & 1167 & 1.34 & 1.31 & 1.29 & 1.29 & 0.016 \\
algebraic & 30\% & both & 4.30--11.4 & 4.29--6.67 & 4.29--4.81 & 4.29--4.67 & 0.001--0.248 \\
algebraic & 20\% & both & 7.30--20.5 & 7.29--11.7 & 7.29--8.30 & 7.29--8.04 & 0.000--0.179 \\
turnover & 60\% & 22.93 & 0 (bound) & $2.3\times10^{7}$ & $2.3\times10^{7}$ & $2.3\times10^{7}$ & 1.961 \\
turnover & 60\% & 1230 & 0 (bound) & $1.7\times10^{7}$ & 12.1 & 16.7 & 2.469 \\
\midrule
\multicolumn{8}{@{}l}{(b) precision with observation noise, $\sdl{\hat\theta}$}\\
\midrule
generator & calibration & $\mathrm{IC}_{50}$ & dose int. & 14.5\,d & monthly & phase-spread & \\
\midrule
linear-clearance (control) & -- & -- & 7.20 & 0.2885 & 0.1893 & 0.1894 & \\
algebraic & 60\% & 63.75 & 4.99 & 0.2861 & 0.1894 & 0.1900 & \\
algebraic & 60\% & 1167 & 8.96 & 0.2912 & 0.1952 & 0.1955 & \\
turnover & 60\% & 22.93 & sing. & sing. & sing. & sing. & \\
turnover & 60\% & 1230 & 355 & sing. & 0.694 & 1.053 & \\
\bottomrule
\end{tabular}

\vspace{1.2ex}
\begin{tabular}{@{}llccccc@{}}
\toprule
\multicolumn{7}{@{}l}{(c) where the load enters the algebraic structure}\\
\midrule
 & & $\theta$ & residual, & residual, & \multicolumn{2}{c}{$\hat\theta/\theta$} \\
\cmidrule(l){6-7}
load enters & marker & identifiable & noiseless & with noise & with priors & without \\
\midrule
not at all (as published) & $E_0+\mathrm{base}\,f$ & no & 1.06, 1.27 & 1.23, 1.59 & 1.41 & inadmissible \\
whole marker & $(E_0+\mathrm{base}\,f)\,u$ & no & 0.18, 0.18 & 0.80, 0.84 & 1.34 & 7.35 \\
suppressible part only & $E_0+\mathrm{base}\,u\,f$ & yes & 0.22, 0.24 & 0.82, 0.88 & 1.35 & 7.73 \\
floor only & $E_0\,u+\mathrm{base}\,f$ & no & 0.95, 1.15 & 1.14, 1.48 & 1.41 & 4.81 \\
separate additive source & $E_0+\mathrm{base}\,f+\gamma u$ & no & 0.93, 1.13 & 1.13, 1.46 & 1.75 & 8.91 \\
\bottomrule
\end{tabular}
% src: audit/PHASE5_ROBUSTNESS.md:522-529,549-556,577-585,1067-1090,1117-1134,1146-1156; audit/k8_recheck.py:221-255
\end{table}

\subsection{The sampling advice transfers}
\label{si:s3:design}

The ranking of schedules is the same under every structure tested. With observation noise, the
monthly and phase-spread designs, the two schedules that cover the dose cycle within 12 visits,
give 0.189--0.196 under the linear-clearance model and under the algebraic structure at both
$\mathrm{IC}_{50}$ roots of the 60\% calibration, agreeing to within about 3\% (Supplementary Table~\ref{stab:structures}b). The
14.5-day schedule, which covers only 43\% of a lap in 12 visits, gives 0.286--0.291 under all
three, and the dose-interval schedule 5.0--9.0. Under the turnover structure the precision is
either unavailable (singular information) or 0.69--1.05 on the monthly and phase-spread schedules
and 355 at the dose interval, so within the usable cells the ranking is unchanged. Results at 2\%
and 10\% day-to-day intake variation differ from these by less than 2\% on the 14.5-day, monthly
and phase-spread schedules.
% src: audit/PHASE5_ROBUSTNESS.md:573-596; results/numbers/lattice.tex:13

The degeneracy behind the aliasing rule also survives. On the dose-interval schedule for a patient
already on therapy, the profile of the fit in $K_m$ stays flat across the whole range scanned
(0.3 to $3\times10^4$\uM) for every generator, so the failure of that schedule is a property of a
record whose draws all share one phase rather than of our parametrisation. Breaking the period
lock bounds $K_m$ under the linear-clearance model and under the algebraic structure, but under
the latter the bounded interval is in the wrong place: 116--165\uM\ against a true 100\uM\ at
the 60\% calibration.
% src: audit/PHASE5_ROBUSTNESS.md:604-635

\subsection{The value of \texorpdfstring{$\hat\theta$}{theta} does not transfer}
\label{si:s3:value}

Fitted to records of the algebraic structure, the linear-clearance model returns $\hat\theta$
1.3--2.0 times the truth at the 60\% calibration (1.29--1.34 on monthly draws), 4.3--11.4 times
at 30\% and 7.3--20.5 times at 20\% (Supplementary Table~\ref{stab:structures}a). $K_m$ moves
with it: at the 30\% calibration the fitted $K_m$ is 462--979\uM\ against a true 100\uM. Yet at
the 60\% calibration the noiseless residual is below a fifth of $\sigma$ on every schedule
(0.000--0.18), and with noise the residual of the monthly and phase-spread fits is
0.78--0.88$\sigma$, against the 0.816 expected of a correct model. The residual does not grow with
the bias: on noiseless monthly records it is 0.016--0.18$\sigma$ at the 60\% calibration,
0.001--0.25$\sigma$ at 30\% and 0.000--0.18$\sigma$ at 20\%, where the bias is largest. Nothing in
the fit signals the error.
% src: audit/PHASE5_ROBUSTNESS.md:522-571

These figures depend on our priors, which shrink the bias but do not remove it. On noiseless
monthly records at the 60\% calibration and 10\% day-to-day intake variation (Supplementary
Table~\ref{stab:structures}c), removing the priors on $k_e$ and on the unrecorded load raises
$\hat\theta/\theta$ from 1.34--1.75 to 4.8--8.9 across the four ways of adding the load
(5.5--8.4 with a second estimator), so the values above are the favourable case. Every
admissible fit lies above the truth, so the direction of the bias does not depend on the
estimator or the priors.
% src: audit/PHASE5_ROBUSTNESS.md:1146-1166; audit/k8_recheck.py:221-255

The bias has a structural cause. With $E=\theta+\Delta$ and $q=1-I_{\max}$,
Eq.~\eqref{seq:algebraic} at a fixed load is a M\"obius function of the increment,
\begin{equation}
\label{seq:moebius}
m(\Delta)=\frac{P\Delta+U}{\Delta+V},\qquad
P=E_0+\mathrm{base}\,q,\quad U=P\theta+(E_0+\mathrm{base})\,\mathrm{IC}_{50},\quad
V=\mathrm{IC}_{50}+\theta ,
\end{equation}
so a record determines three numbers $(P,U,V)$, while the structure has four free parameters at
fixed $I_{\max}$. For any $\theta$ between 0 and $V$, setting $\mathrm{IC}_{50}=V-\theta$,
$\mathrm{base}=(U-PV)/\bigl(I_{\max}(V-\theta)\bigr)$ and $E_0=P-q\,\mathrm{base}$ reproduces the
record exactly. Walking this family, $\theta$ runs from 3.06 to 3033\um, a factor of 990, while the
record changes by at most $1.4\times10^{-14}$ in relative terms, with or without variation in the
load. Under this structure $\theta$ is therefore not identifiable, at any design and any number
of visits, and our estimator reports one point of the family, chosen by our parametrisation and
priors. The same count applies to the linear-clearance model: its quasi-steady map,
Eq.~\eqref{seq:s1-invariants}, carries three numbers for four parameters at a constant load, and
day-to-day variation in the recorded load supplies the fourth (Section~\ref{si:s1-si}). In the algebraic structure, with the load scaling the whole marker,
variation in the recorded load determines $b$ but adds nothing to $(P,U,V)$, and the structure
has one more parameter than the linear-clearance model: a free marker floor $E_0$ alongside a
free amplitude base, which enter the record only through $P$ and $E_0+\mathrm{base}$. It
therefore remains one number short. Fixing $E_0=0$ or measuring $\mathrm{IC}_{50}$
independently makes $\theta$ identifiable again.
% src: paper/sections/A1_supplementary.tex:84-101,115-123; audit/PHASE5_ROBUSTNESS.md:828-857,893-913,939-1014,1194-1230

\subsection{Where the load enters}
\label{si:s3:load}

Because the algebraic structure has no load, whether $\theta$ is identifiable under it turns on a
choice its source does not make. Supplementary Table~\ref{stab:structures}c lists four plausible
ways of adding the load, with the structure as published (no load) for reference. At a varying
recorded load, three of the four leave $\theta$ unidentifiable (the whole marker scaled by the
load, the load lifting only the floor $E_0$, or the load as a separate additive source), and one
restores it: a load that lifts only the suppressible part base. On the definitions of the original
model (base, the marker amenable to suppression; $E_0$, the level not affected by the therapy) the
last is arguably the most natural reading, but that is our judgement. At a constant load every
choice leaves $\theta$ unidentifiable.
% src: audit/PHASE5_ROBUSTNESS.md:1067-1115; paper/sections/A1_supplementary.tex:103-113

The bias is above 1 for every choice (1.34--1.75 on noiseless monthly records,
Supplementary Table~\ref{stab:structures}c; 1.32--1.76 on medians of noisy monthly and
phase-spread fits), but whether the fit reveals it depends on the choice. The residual stays near
or below a fifth of $\sigma$ when the load scales the whole marker (0.18) or only its suppressible part (0.22--0.24), and with noise it is
then indistinguishable from the 0.816 of a correct model (0.80--0.88). When the load enters
through the floor, as an additive source, or not at all, the noiseless residual is about one
$\sigma$ (0.93--1.27) and the residual with noise 1.13--1.59, which an analyst could detect as an
excess over the expected value.
% src: audit/PHASE5_ROBUSTNESS.md:1117-1144,1152-1156

\subsection{The turnover structure}
\label{si:s3:turnover}

Under the turnover structure $\theta$ is identifiable at a constant load with no outside
information. Its forcing term in Eq.~\eqref{seq:turnover} is again a ratio of first-degree
functions of the increment, $(q\Delta+W)/(\Delta+V)$ with $W=q\theta+\mathrm{IC}_{50}$ and
$V=\theta+\mathrm{IC}_{50}$, but with no free scale in front of it, since $k_{\mathrm{in}}$ is the
recorded load; so $W$ and $V$ give $\theta=(V-W)/I_{\max}$, and because the marker is a state its
transients fix $k_{\mathrm{out}}$ separately. Moving $\theta$ away from the truth while holding $V$
and re-optimising $k_{\mathrm{out}}$, the nearest admissible alternative tested (2970\um) leaves
a mismatch $2.8\times10^9$ times that at the truth. The marker state supplies the fourth number,
and the structure has no free marker scale that would cancel it.
% src: audit/PHASE5_ROBUSTNESS.md:844-857,915-933,976-1000

The linear-clearance model cannot fit records of this structure, and the misfit shows on the
recommended schedules. On the 14.5-day, monthly and phase-spread schedules the noiseless residual
is 2.0--2.6$\sigma$ and $\hat\theta$ runs away, to 12 to $2.3\times10^{7}$ times the truth,
into the limit where the enzymatic term becomes indistinguishable from the linear route. On the
dose-interval schedules the residual is small (0.13--0.62$\sigma$ from the start of therapy, 0.000
for a patient already on therapy) and $\hat\theta$ goes to the estimator's lower bound instead.
At the dose interval the misfit therefore goes unseen and the fit returns no capacity, whereas on
draws across the cycle it is visible.
% src: audit/PHASE5_ROBUSTNESS.md:522-547

\subsection{One draw after an infusion}
\label{si:s3:draw}

The linear-clearance model and the algebraic structure can share the whole steady map and still
differ in time. In the linear-clearance model the marker is a state that relaxes towards its
new level with the time constant of Section~\ref{si:s2:firsthour}, about 0.8\,h before an
infusion and shorter after it; in the algebraic structure the marker follows the enzyme instantly. Just after an
infusion the delivered enzyme jumps to its peak: the algebraic marker drops at once to its new
level, while the linear-clearance marker has barely moved. To isolate this difference we
calibrated an algebraic marker to agree with the linear-clearance model at both ends of the
transient, the trough before the infusion and the settled level after it, so that the two share the steady map exactly and differ only in how they
travel between levels. With a 5\% assay coefficient of variation applied at each draw's level, one
draw separates them by 16.3 standard deviations 3 minutes after the infusion, 9.7 at 15 minutes,
4.4 at 30 minutes, 0.68 at one hour and 0.01 at two hours. A single draw gives at least three
standard deviations at any assay precision up to 20\%; at a deliberately poor 40\% it takes five
draws spread over the first hour. These values are for the linear-clearance model at normal
capacity.
% src: results/numbers/structure.tex:4-9; paper/sections/08_specification.tex:71-100; audit/t7_structure_transient.py:118-215

The window is, if anything, earlier in the carnitine model. Over one hour carnitine moves only
2.9\% of the way to its new value, so during the transient the conjugation route acts at the
pre-infusion carnitine level, and the relaxation time at matched capacity is 0.555\,h (0.587\,h
with carnitine eliminated at equilibrium) against 0.81\,h in the linear-clearance model. Draw
times designed on the linear-clearance model would therefore sample late; draws in the first 15
minutes carry almost all of the information. Such a draw is excluded from the steady-state fit
(Section~\ref{si:s2:firsthour}) but should be part of the protocol.
% src: results/numbers/structure.tex:3,10-11; results/numbers/design.tex:110; paper/sections/08_specification.tex:93-100

The draw rules out a family of structures rather than identifying one. It rejects the algebraic
structure, under which, as published, $\theta$ is not identifiable from a marker record at any
design, and any other structure in which the marker has no relaxation of its own. It does not
separate the linear-clearance model from a structure whose marker is also a state but whose
steady map differs. We did not compute the separation for the turnover structure, against
which the routine record itself already shows the misfit
(Section~\ref{si:s3:turnover}).
% src: paper/sections/08_specification.tex:84-91; paper/sections/11_scope.tex:295-306; audit/t7_structure_transient.py:166-186

\subsection{What follows}
\label{si:s3:follows}

Across the structures tested, the design rules hold: sample across the dose cycle, and never at
the dose interval. Under the linear-clearance model and the algebraic structure the monthly and
phase-spread schedules also give the same precision to within about 3\%. A value of $\hat\theta$ depends on the assumed model structure as well as on the patient.
Under the algebraic structure as published, and in three of the four ways of adding the load to
it, $\theta$ is not identifiable; a fit of the linear-clearance model to its records returns
$\hat\theta$ 1.3 to 20.5 times too high with priors and several times higher without them; and
when the load scales the whole marker or only its suppressible part, the residual gives no
warning. Because a wrong structure can fit the routine record well, the check on structure has
to come from a measurement outside that record. The draw about 3 minutes after an infusion
(Tier~1 of the main text's data specification) is that check for the algebraic family. The same
caveat applies to the $K_m$ the record recovers (Section~\ref{si:s2:record}).
% src: paper/sections/A1_supplementary.tex:125-136; paper/sections/11_scope.tex:295-306; audit/PHASE5_ROBUSTNESS.md:522-602
